\documentclass[a4paper,11pt]{letter}
\usepackage[default]{sourcesanspro}
\usepackage{graphicx} % Required for inserting images
\usepackage{geometry}
\usepackage{fancyhdr}
\usepackage{fontawesome}  % Font awesome icons
\usepackage[utf8]{inputenc}
\usepackage{setspace}
\usepackage[hidelinks]{hyperref}


% Set the page margins
\geometry{a4paper, top=4cm, bottom=3cm, left=2.5cm, right=2.5cm}

% Setup the header
\pagestyle{fancy}
\fancyhf{} % clear all header fields
\fancyhead[L]{
    {\large \textbf{Julio Molleda}}\\
    \href{mailto:jmolleda@uniovi.es} jmolleda@uniovi.es\\
    University of Oviedo, Dept. Comp. Sci. and Eng., Gijón, Spain \\
    
    \medskip
    Computers and Electronics in Agriculture\\
    \textit{Foundation Models in Agriculture: Advancing AI for Precision and Smart Farming} \\
    \medskip
    }
\fancyhead[R]{\textbf{\today}}
\medskip\vspace*{35pt} 

%%%%%%  LETTER STARTS HERE  %%%%%%%%%%%%%%%%%%%%%%%%%%%%

\begin{document}

%Salutation
\normalsize Dear Editors,

% Letter body
\small \noindent
We wish to submit an original research paper entitled “Evaluation Framework for Multimodal Language Models Applied to Geospatial Open Data in Agriculture” to be published by the journal \textit{Computers and Electronics in Agriculture} for the special issue "Foundation Models in Agriculture: Advancing AI for Precision and Smart Farming". 

\vspace{1pt}
This paper addresses a critical gap in the systematic evaluation of multimodal language models augmented with Earth observation multi-domain data for precision agriculture and land management. Existing evaluation frameworks fail to measure how these models actually utilize dynamically retrieved context. To solve this, we propose a novel, reproducible evaluation methodology that quantifies the marginal contributions of distinct architectural modules by applying collaborative game theory analysis. We showcase this framework by applying it to a previously validated context-amplified LLM, yielding very significant findings on both the system architecture and the methodological framework itself. 

\vspace{1pt}
Our approach introduces a hybrid human-in-the-loop panel combining two human evaluators and two distinct LLM judges, to assess a 143-query dataset across five targeted metrics. By designing a full factorial ablation study modeled as a cooperative game and applying Shapley values, we quantify the impact of Satellite, Orthophoto, and Retrieval-Augmented Generation (RAG) modules. We show that visual data drive up to a 32$\%$ increase in overall response utility for given queries, while RAG module exhibits high, query-dependent volatility. Furthermore, our context length analysis proves the architecture's resilience to information overload, while inter-evaluator consensus metrics highlight the continued necessity of human oversight to adequately penalize bias.


\vspace{1pt}
To ensure reproducibility and open science, the complete evaluation dataset, containing agricultural parcel orthophotos, retrieved external context, and generated responses, is shared supplementary material.


\vspace{1pt}
We believe this study aligns perfectly with the scope of \textit{Computer and Electronics in Agriculture}, and would be a particularly strong fit for the special issue "Foundation Models in Agriculture: Advancing AI for Precision and Smart Farming". We feel that our findings will be highly relevant to your audience, providing a solid contribution to the research community through the establishment of a rigorous, mathematical standard for benchmarking next-generation foundation models in smart farming.

\vspace{1pt}
The article “Evaluation Framework for Multimodal Language Models Applied to Geospatial Open Data in Agriculture” has not been published elsewhere, and it reflects original research conducted by its authors. None of the authors have any conflicts of interest to disclose concerning this study. 


Thank you for your time and consideration. 

\bigskip
\noindent
Sincerely,\\
\textbf{Julio Molleda}

\end{document}